\documentclass[11pt,letterpaper]{article}
\usepackage{fontspec}
\setmainfont{Latin Modern Roman}
\usepackage{amsmath,amssymb,amsthm,mathtools}
\usepackage[margin=1.05in]{geometry}
\usepackage{microtype}
\usepackage{xcolor}
\definecolor{paperlink}{RGB}{0,0,114}
\usepackage{hyperref}
\hypersetup{colorlinks=true,linktoc=all,linkcolor=paperlink,citecolor=paperlink,urlcolor=paperlink,pdfstartview=Fit,pdftitle={A Compact-Support Obstruction in Wang's Planar Strong-Solution Class}}
\usepackage{fancyhdr}
\pagestyle{fancy}
\fancyhf{}
\fancyhead[L]{\small Compact-support obstruction for planar full flow}
\fancyhead[R]{\small\thepage}
\setlength{\headheight}{14pt}
\setlength{\parindent}{1.25em}
\setlength{\parskip}{0pt}
\setlength{\emergencystretch}{1em}
\numberwithin{equation}{section}
\newtheorem{theorem}{Theorem}[section]
\newtheorem{lemma}[theorem]{Lemma}
\newtheorem{proposition}[theorem]{Proposition}
\newtheorem{corollary}[theorem]{Corollary}
\theoremstyle{definition}
\newtheorem{definition}[theorem]{Definition}
\newcommand{\keywords}[1]{\par\noindent\textbf{Keywords.} #1\par}
\newcommand{\R}{\mathbb R}
\newcommand{\dd}{\,dx}
\newcommand{\dt}{\,dt}
\newcommand{\Div}{\operatorname{div}}
\newcommand{\supp}{\operatorname{supp}}
\newcommand{\norm}[2]{\lVert#1\rVert_{#2}}
\newcommand{\cH}{\mathcal H}
\title{A Compact-Support Obstruction in Wang's Planar\\
Full Compressible Navier--Stokes Strong-Solution Class}
\author{}
\date{}
\begin{document}
\maketitle

\begin{abstract}
We consider the planar full compressible Navier--Stokes system in
Wang's version 3 strong-solution class, with positive heat conductivity
and nonnegative temperature. In the strong
solution class displayed in Wang's 2023 preprint, we prove a conditional
obstruction to global solutions with compactly supported initial
density, positive mass, and positive physical energy, under the strict
viscosity condition $\mu+\lambda>0$. Finite spatial Lebesgue norms force
the temperature and velocity to vanish in the transported exterior
vacuum. A common set of admissible times then fixes the initial support
ball. Compact support supplies the time regularity needed to recover the
initial kinetic energy from the conservative momentum trace; the
thermal trace is handled directly by the conservative equation.
Conservation of total energy and the virial identity give an explicit
upper bound for every interval on which the stated class is maintained.
A smooth compatible family has fixed positive mass, common bounds for
all the stipulated initial norms, and positive energy tending to zero.
The theorem excludes global solutions retaining the displayed class;
it does not independently reconstruct the cited local solution.
\end{abstract}
\noindent\textbf{Manuscript status.} Analytic proof draft; the exact PDE theorem has not been fully formalized in Lean.\par\medskip
\keywords{Compressible Navier--Stokes equations $\cdot$
Heat conduction $\cdot$ Planar vacuum $\cdot$
Compact support $\cdot$ Physical energy}

\section{The equations and the conditional obstruction}\label{sec:intro}

We study the Cauchy problem for the heat-conducting compressible
Navier--Stokes system on $\R^2$:
\begin{align}
 \rho_t+\Div(\rho u)&=0,\label{eq:mass}\\
 (\rho u)_t+\Div(\rho u\otimes u)-\Div S+\nabla P&=0,
 \label{eq:momentum}\\
 c_v\bigl[(\rho\theta)_t+\Div(\rho\theta u)\bigr]
 &=Q(u)-P\Div u+\kappa\Delta\theta.\label{eq:thermal}
\end{align}
Here $\rho$, $u=(u_1,u_2)$, and $\theta$ denote the density, velocity,
and absolute temperature. The pressure and viscous stress are
\begin{equation}
 \begin{aligned}
 P&=R\rho\theta,\qquad
 D(u)=\tfrac12(\nabla u+\nabla u^{\mathsf T}),\\
 S&=2\mu D(u)+\lambda(\Div u)I,\qquad
 Q(u)=S:\nabla u
       =2\mu|D(u)|^2+\lambda(\Div u)^2.
 \end{aligned}\label{eq:stress}
\end{equation}
We assume throughout that
\begin{equation}
 \mu>0,\qquad b:=\mu+\lambda>0,\qquad R,c_v,\kappa>0.
 \label{eq:parameters}
\end{equation}
The source condition \cite[(1.2)]{WangV3} permits $b\ge0$;
our conclusion is restricted to $b>0$.
In particular, $\Div S=\mu\Delta u+b\nabla\Div u$. The letter $R$
always denotes the gas constant; support balls will have radius $L$.

The argument concerns a precise solution class. In
\cite[Theorem 1.1]{WangV3}, Wang states local existence for this system
with vacuum and gives a complete list of weighted and unweighted
regularity properties. We use that displayed list as a hypothesis on
a putative solution and do not invoke the local existence assertion
to construct one. Thus the conclusion below excludes global
solutions in the stated class. Finite-time breakdown of a particular
locally constructed solution additionally requires a separate local
construction.

The exterior thermal equation is decisive. In a spacetime vacuum
region it becomes $-\kappa\Delta\theta=Q(u)$. In two dimensions,
nonnegativity and finite spatial Lebesgue integrability force such a
superharmonic temperature to vanish on a connected exterior domain.
Under \eqref{eq:parameters}, the resulting equality $Q(u)=0$ makes the
velocity a rigid motion there, and its finite Lebesgue norm eliminates
that motion. We first carry out this argument on the transported
exterior of the initial support ball. After proving that the exterior
flow is stationary, we use the resulting fixed support bound in the
energy and virial identities.

All function spaces and spatial integrals without a specified domain
are over $\R^2$. On a fixed finite time interval, constants in
intermediate estimates may depend on that interval, the physical
parameters, the exponents, the support radius, and the finite solution
norms explicitly used. The constant in the final lifespan bound is
given exactly.

Fix $a>1$, $q>2$, and $\eta_0>0$, and set
\begin{equation}
 \begin{aligned}
 \bar x&=(e+|x|^2)^{1/2}\log^{1+\eta_0}(e+|x|^2),\\
 \alpha&=\min\{a/2,1\},\qquad
 p_u=\frac4\alpha,\qquad
 p_\theta=\frac6{2\alpha-1}.
 \end{aligned}\label{eq:weights}
\end{equation}
Thus $1/2<\alpha\le1$, $4\le p_u<8$, and
$6\le p_\theta<\infty$. The exponents $p_u,p_\theta$ are denoted
by $q_1,q_2$ in \cite[(1.8)]{WangV3}.

The initial assumptions are the displayed data conditions
\cite[(1.5)--(1.6)]{WangV3}:
\begin{equation}
 \begin{aligned}
 &\rho_0,\theta_0\ge0,\qquad
       \bar x^a\rho_0\in L^1\cap H^1\cap W^{1,q},\\
 &\sqrt{\rho_0}u_0,\ \sqrt{\rho_0}\theta_0\in L^2,\qquad
       \nabla u_0\in H^1,\quad \nabla\theta_0\in L^2,
 \end{aligned}\label{eq:data}
\end{equation}
and
\begin{equation}
 -\mu\Delta u_0-b\nabla\Div u_0
       +R\nabla(\rho_0\theta_0)=\sqrt{\rho_0}\,g,
 \qquad g\in L^2.
 \label{eq:compatibility}
\end{equation}
In particular, $\norm{\sqrt{\rho_0}\theta_0}{2}^2
=\int\rho_0\theta_0^2\dd$. No unweighted $L^2$ assumption on
$u_0$ or $\theta_0$, finite-$L^{p_u}$ assumption on $u_0$, or thermal
compatibility condition is included in these displayed data conditions.

\begin{definition}\label{def:solution}
A solution in the specified class on $[0,T_*)$, where $T_*>0$, consists
of $\rho,u,\theta$ with $\rho,\theta\ge0$, for which all derivatives
in \eqref{eq:mass}--\eqref{eq:thermal} are regular distributions and
the equations hold almost everywhere. For every $0<T<T_*$ it satisfies
the complete regularity list
\begin{equation}
 \begin{aligned}
 &\rho\in C([0,T];L^1\cap H^1\cap W^{1,q}),\qquad
   \bar x^a\rho\in L^\infty(0,T;L^1\cap H^1\cap W^{1,q}),\\
 &\nabla u\in L^\infty(0,T;H^1)\cap L^2(0,T;W^{1,q}),
   \qquad u\in L^\infty(0,T;L^{p_u}),\\
 &\nabla\theta,\ \sqrt t\,\nabla^2\theta
     \in L^\infty(0,T;L^2)\cap L^2(0,T;H^1),
   \qquad\theta\in L^2(0,T;L^{p_\theta}),\\
 &\sqrt\rho\,u,\ \sqrt\rho\,\theta,\ \sqrt\rho\,u_t,\
    \sqrt t\sqrt\rho\,\theta_t,\
    \bar x^{\alpha/2}\nabla u\in L^\infty(0,T;L^2),\\
 &\nabla u_t,\ \sqrt\rho\,\theta_t,\
     \sqrt t\,\nabla\theta_t\in L^2(\R^2\times(0,T)).
 \end{aligned}\label{eq:fullclass}
\end{equation}
Its initial conditions are imposed on conservative variables:
\begin{equation}
 \begin{gathered}
 \rho(0)=\rho_0,\qquad
 \rho u(t)\longrightarrow\rho_0u_0,\qquad
 \rho\theta(t)\longrightarrow\rho_0\theta_0\\
 \hbox{in }\mathcal D'(\R^2)\quad(t\downarrow0).
 \end{gathered}
 \label{eq:traces}
\end{equation}
The density trace is also strong in the spaces specified in
\eqref{eq:fullclass}.
\end{definition}

The list \eqref{eq:fullclass} reproduces \cite[(1.7)]{WangV3};
the interpretation of strong derivatives is
\cite[Definition 1.1]{WangV3}. The conservative initial variables
are those of \cite[(1.3)]{WangV3}; we interpret (1.3) as distributional
convergence of those variables, without imposing an unweighted initial trace.
The far-field condition
\cite[(1.4)]{WangV3} is qualified by \emph{in some weak sense}.
We do not assign an additional topology to that phrase. Our proof
uses the explicit finite norms of $u$ and $\theta$ in
\eqref{eq:fullclass}, imposed on the functions themselves.
Replacing them by conditions on representatives modulo constants
would change the hypothesis.

For clarity, the following portion of the full class suffices
for the proof:
\begin{equation}
 \begin{aligned}
 &\rho\in C([0,T];L^1\cap H^1\cap W^{1,q}),\\
 &\nabla u\in L^\infty(0,T;H^1)
                    \cap L^2(0,T;W^{1,q}),\\
 &u\in L^\infty(0,T;L^{p_u}),\\
 &\sqrt\rho\,u,\ \sqrt\rho\,\theta,\ \nabla\theta
                       \in L^\infty(0,T;L^2),\\
 &\nabla u_t,\ \nabla^2\theta\in L^2(0,T;L^2),\qquad
       \theta\in L^2(0,T;L^{p_\theta}).
 \end{aligned}\label{eq:usedclass}
\end{equation}
This identifies the analytic inputs; Definition~\ref{def:solution}
retains the full original list. The unweighted spacetime Hessian
bound in \eqref{eq:usedclass} is an explicit source assumption,
not a consequence asserted from its time-weighted counterpart.

Here and below $B_L=B_L(0)$ is the origin-centered ball. Suppose now that,
for some $L>0$,
\begin{equation}
 \begin{gathered}
 \supp\rho_0\subset\overline{B_L},\qquad M:=\int\rho_0\dd>0,\\
 E_0:=\int\left(\tfrac12\rho_0|u_0|^2+c_v\rho_0\theta_0\right)\dd>0.
 \end{gathered}
 \label{eq:compactdata}
\end{equation}
Positive mass is explicit here; no mass normalization is used.
The energy is finite by \eqref{eq:data}, since
$\int\rho_0\theta_0\dd
\le M^{1/2}\norm{\sqrt{\rho_0}\theta_0}{2}$.
Define
\begin{equation}
 I_0=\int\rho_0|x|^2\dd,\qquad
 J_0=\int\rho_0u_0\cdot x\dd,\qquad
 \beta=\min\{1,R/c_v\}>0.
 \label{eq:moments}
\end{equation}
Both moments are finite:
$I_0\le ML^2$ and
$|J_0|\le I_0^{1/2}\norm{\sqrt{\rho_0}u_0}{2}$. In fact,
\begin{equation}
 I_0<ML^2,
 \label{eq:strictmoment}
\end{equation}
because $\rho_0$ is an integrable function of positive mass and the
circle $\{|x|=L\}$ has two-dimensional Lebesgue measure zero.

\begin{theorem}\label{thm:obstruction}
Assume \eqref{eq:parameters}, \eqref{eq:data},
\eqref{eq:compatibility}, and \eqref{eq:compactdata}.
If a solution in the sense of Definition~\ref{def:solution} exists on
$[0,T_*)$, then for every $0\le t<T_*$,
\begin{equation}
 I_0+2J_0t+2\beta E_0t^2
       \le\int\rho(t,x)|x|^2\dd\le ML^2.
 \label{eq:mainbound}
\end{equation}
Consequently,
\begin{equation}
 T_*\le
 \frac{-J_0+\sqrt{J_0^2+2\beta E_0(ML^2-I_0)}}
      {2\beta E_0}<\infty.
 \label{eq:lifespan}
\end{equation}
In particular, no global solution satisfying the stated class on
every finite interval exists for these data.
\end{theorem}

The theorem requires neither nonzero total momentum nor a positive
density or temperature lower bound. Its second moment follows from
compact support. If a local solution has been constructed
independently, \eqref{eq:lifespan} limits any extension maintaining
this class. Without that construction, the conclusion remains the
conditional absence of a global solution.

\section{Exterior vacuum and fixed support}\label{sec:exterior}

We first prove the planar rigidity fact used for the temperature.

\begin{lemma}\label{lem:superharmonic}
Let $\Omega\subset\R^2$ be connected and open, and suppose that
$\{|x|>L_1\}\subset\Omega$ for some $L_1>0$. If
\[
 v\in W^{2,2}_{\mathrm{loc}}(\Omega)\cap L^p(\Omega),\qquad
 v\ge0,\qquad -\Delta v\ge0
\]
in distributions, where $1\le p<\infty$, then $v=0$ in $\Omega$.
\end{lemma}

\begin{proof}
Local Sobolev regularity gives a continuous representative of $v$.
For $s>\log L_1$, let
\[
 h(s)=\frac1{2\pi}\int_0^{2\pi}
           v(e^s\cos\varphi,e^s\sin\varphi)\,d\varphi.
\]
The polar change of variables on compact annuli and angular
integration show that $h\in H^2_{\mathrm{loc}}$ and
\begin{equation}
 h''(s)=\frac{e^{2s}}{2\pi}\int_0^{2\pi}
             \Delta v(e^s\cos\varphi,e^s\sin\varphi)\,d\varphi
       \le0
 \label{eq:concavemean}
\end{equation}
almost everywhere. Thus $h$ is nonnegative and concave on a half-line.
It is nondecreasing: a negative secant slope, together with concavity,
would force $h$ to become negative at sufficiently large $s$.
Jensen's inequality gives
\begin{equation}
 2\pi\int_{\log L_1}^{\infty} h(s)^p e^{2s}\,ds
       \le\int_{|x|>L_1}|v(x)|^p\dd<\infty.
 \label{eq:meanLp}
\end{equation}
If $h(s_0)>0$, monotonicity would make the left side infinite.
Hence $h=0$, and nonnegativity and continuity imply $v=0$ on
$\{|x|>L_1\}$.

To propagate this equality, the local super-mean inequality gives
\begin{equation}
 \frac1{2\pi}\int_0^{2\pi}
       v(z+r(\cos\varphi,\sin\varphi))\,d\varphi\le v(z),
 \qquad \overline{B_r(z)}\subset\Omega.
 \label{eq:supermean}
\end{equation}
For the stated regularity, convolve $v$ inside a slightly larger
ball, apply the smooth inequality, and pass to the limit locally
uniformly. If $v(z)=0$, nonnegativity forces the average in
\eqref{eq:supermean} to vanish for every sufficiently small $r$.
Continuity therefore gives $v=0$ in a neighborhood of $z$.
The zero set is relatively open and closed in $\Omega$ and is
nonempty. Connectedness proves the assertion.
\end{proof}

\begin{proposition}\label{prop:fixedsupport}
Let a solution satisfy Definition~\ref{def:solution} on $[0,T_*)$,
and suppose that $\supp\rho_0\subset\overline{B_L}$. Then
\begin{equation}
 \begin{gathered}
 \supp\rho(t)\subset\overline{B_L}\quad(0\le t<T_*),\\
 u(t)=\theta(t)=0\quad\hbox{on }\R^2\setminus\overline{B_L}
 \end{gathered}
 \label{eq:fixedsupport}
\end{equation}
for almost every $t\in(0,T_*)$. The density assertion holds for its
continuous spatial representative.
\end{proposition}

\begin{proof}
Fix $0<T<T_*$. Since $q>2$, Sobolev embedding applied to $\nabla u$
gives
$\norm{\nabla u(t)}{\infty}
\le C_q\norm{\nabla u(t)}{W^{1,q}}$
for almost every $t$. Averaging the velocity over a unit ball and
using its Lipschitz bound, we obtain
\begin{align}
 |u(t,x)|
 &\le \frac1{|B_1|}\int_{B_1(x)}|u(t,y)|\,dy
           +C\norm{\nabla u(t)}{\infty}\notag\\
 &\le C_{p_u}\norm{u(t)}{p_u}
           +C\norm{\nabla u(t)}{\infty}.
 \label{eq:velocitybound}
\end{align}
The source norms and Cauchy--Schwarz in time imply
\begin{equation}
 A_T:=\int_0^T\norm{u(t)}{\infty}\dt<\infty,\qquad
 B_T:=\int_0^T\norm{\nabla u(t)}{\infty}\dt<\infty.
 \label{eq:flowbounds}
\end{equation}
The spatial representative is $C^1$ for almost every time, since
$\nabla u\in W^{1,q}$. Values on the remaining null set of times
do not affect the integral flow
\begin{equation}
 X(t,y)=y+\int_0^t u(s,X(s,y))\,ds.
 \label{eq:flow}
\end{equation}
Existence and forward and backward uniqueness follow from
\eqref{eq:flowbounds}. Gronwall's inequality gives
\begin{equation}
 \begin{aligned}
 |X(t,y)-y|&\le A_T,\\
 e^{-B_T}|y-z|&\le |X(t,y)-X(t,z)|
                        \le e^{B_T}|y-z|.
 \end{aligned}\label{eq:bilipschitz}
\end{equation}
Thus $X(t,\cdot)$ is a bi-Lipschitz homeomorphism. The flow and its
inverse are jointly continuous in space and time, and
\begin{equation}
 \mathcal J(t,y)=\det\nabla_yX(t,y)
    =\exp\left(\int_0^t\Div u(s,X(s,y))\,ds\right)>0.
 \label{eq:jacobian}
\end{equation}
The variational equation for spatial derivatives is justified by
the spatial $C^1$ representative and the integrable derivative bound;
it also follows by spatial smoothing and Lipschitz stability of the
flow. In particular, $\mathcal J$ has positive upper and lower bounds
depending only on $B_T$.

We justify the density characteristic identity rather than assume
a pointwise time chain rule. For a spatial mollifier $\zeta_\delta$
put $\rho_\delta=\zeta_\delta*\rho$. The continuity equation yields
\begin{equation}
 (\rho_\delta)_t+u\cdot\nabla\rho_\delta
           +(\Div u)\rho_\delta=r_\delta.
 \label{eq:mollifieddensity}
\end{equation}
Its transport commutator satisfies
\begin{equation}
 \norm{u\cdot\nabla\rho_\delta
               -(u\cdot\nabla\rho)_\delta}{q}
       \le C\delta\norm{\nabla u}{\infty}\norm{\nabla\rho}{q}.
 \label{eq:transportcommutator}
\end{equation}
The other commutator,
$(\Div u)\rho_\delta-((\Div u)\rho)_\delta$, tends to zero in $L^q$
for almost every time and is bounded by
$2\norm{\Div u}{\infty}\norm{\rho}{q}$.
Hence $r_\delta\to0$ in $L^1(0,T;L^q)$.
Integrate \eqref{eq:mollifieddensity} along \eqref{eq:flow}.
The Jacobian bounds allow passage to the limit in $L^q_y$, giving
\begin{equation}
 \rho(t,X(t,y))\mathcal J(t,y)=\rho_0(y).
 \label{eq:transportdensity}
\end{equation}
Initially this holds almost everywhere in $y$. Spatial continuity
of $\rho$, the flow, and \eqref{eq:jacobian} extends it to every $y$.
Continuity of $\rho$ in time in $W^{1,q}$ extends it to every
$t\in[0,T]$. No division by density is involved.

Put $\Omega_0=\R^2\setminus\overline{B_L}$ and
$\Omega_t=X(t,\Omega_0)$. Equation \eqref{eq:transportdensity} gives
$\rho(t)=0$ on $\Omega_t$ for every $t$. Each $\Omega_t$ is connected
and open, with compact complement $X(t,\overline{B_L})$, and
\begin{equation}
 \{|x|>L+A_T+1\}\subset\Omega_t,\qquad 0\le t\le T.
 \label{eq:commonexterior}
\end{equation}
At this stage the support is only transported.
The spacetime set
\[
 \mathcal V=\{(t,x):0<t<T,\ X(t,\cdot)^{-1}(x)\in\Omega_0\}
\]
is open by joint continuity of the inverse flow. On $\mathcal V$,
$\rho=0$, so $\rho\theta$, $\rho\theta u$, and $P$ vanish as
distributions. Restriction of the conservative thermal equation gives
\begin{equation}
 -\kappa\Delta\theta=Q(u)\quad\hbox{in }\mathcal D'(\mathcal V).
 \label{eq:vacuumthermal}
\end{equation}
This uses an open spacetime vacuum region; it does not set a time
derivative to zero merely from vacuum on a single time slice.

To slice \eqref{eq:vacuumthermal}, cover $\mathcal V$ by countably
many spacetime cylinders with closures contained in $\mathcal V$.
For a countable dense collection of spatial tests in each cylinder,
Fubini's theorem removes a single null set of times. Taking the
countable union of these exceptional sets, together with the null
sets for the source norms and signs, gives one full-measure set
$G\subset(0,T)$ such that, for each $t\in G$,
\begin{equation}
 \begin{gathered}
 \theta(t)\in W^{2,2}_{\mathrm{loc}}(\R^2)\cap L^{p_\theta},
       \qquad \theta(t)\ge0,\\
 u(t)\in L^{p_u},\qquad \nabla u(t)\in W^{1,q},\\
 -\kappa\Delta\theta(t)=Q(u(t))
       \quad\hbox{in }\mathcal D'(\Omega_t).
 \end{gathered}\label{eq:commontimes}
\end{equation}
Local $L^2$ integrability of $\theta(t)$ follows from $p_\theta\ge2$;
the gradient and unweighted Hessian bounds provide the other local
derivatives. Thus $\theta(t)$ is continuous. Crucially, $G$ has been
chosen independently of all trajectories and spatial points.

Write $d=\Div u$ and $D_0=D(u)-\tfrac12dI$. In two dimensions,
\begin{equation}
 Q(u)=2\mu|D_0|^2+b\,d^2\ge0.
 \label{eq:positiveQ}
\end{equation}
For every $t\in G$, Lemma~\ref{lem:superharmonic} applies on
$\Omega_t$ with $p=p_\theta$. It gives $\theta(t)=0$ there.
Then \eqref{eq:vacuumthermal} gives $Q(u(t))=0$ there, and strict
$\mu>0$, $b>0$ imply $D(u(t))=0$. The distributional identity
\begin{equation}
 \partial_j\partial_k u_i
   =\partial_j D_{ik}(u)+\partial_k D_{ij}(u)
                         -\partial_i D_{jk}(u)
 \label{eq:rigididentity}
\end{equation}
shows that all second derivatives of $u(t)$ vanish on $\Omega_t$.
Connectedness gives $u(t,x)=A(t)x+c(t)$, where $A(t)$ is a constant
skew-symmetric matrix. A nonzero constant or planar rotation has
infinite $L^{p_u}$ norm on a full exterior region. Hence
\begin{equation}
 u(t,x)=\theta(t,x)=0
       \quad(x\in\Omega_t,\ t\in G).
 \label{eq:exteriorzero}
\end{equation}
Spatial continuity makes these equalities pointwise.

Fix any $y\in\Omega_0$. Since $X(t,y)\in\Omega_t$,
\eqref{eq:exteriorzero} gives $u(t,X(t,y))=0$ on the same set $G$
for every such $y$. Absolute continuity of the trajectories yields
\[
 X(t,y)=y\quad(0\le t\le T,\ y\in\Omega_0).
\]
Continuity fixes $\partial B_L$ as well. Bijectivity therefore implies
$X(t,\overline{B_L})=\overline{B_L}$ and $\Omega_t=\Omega_0$.
Equations \eqref{eq:transportdensity} and \eqref{eq:exteriorzero}
now prove \eqref{eq:fixedsupport}. The argument applies to every
$T<T_*$.
\end{proof}

The fixed initial ball is essential: a radius growing with
$\int_0^t\norm{u}{\infty}$ would not give the uniform moment bound
in \eqref{eq:mainbound}. Strict viscosity is also essential to
this proof. At $b=0$, \eqref{eq:positiveQ} does not control
$\Div u$, so $Q(u)=0$ does not force $D(u)=0$. No endpoint
assertion is made here.

\section{Time traces and conservation of energy}\label{sec:energy}

We now justify the total-energy constant, including its value at
$t=0$. Conservative momentum convergence alone does not imply
convergence of the kinetic energy. The additional time regularity
required for that step follows from Proposition~\ref{prop:fixedsupport}.

\begin{proposition}\label{prop:energy}
Under the hypotheses of Theorem~\ref{thm:obstruction}, set
\[
 K(t)=\frac12\int\rho|u|^2\dd,\qquad
 U(t)=c_v\int\rho\theta\dd.
\]
For each $T<T_*$ these quantities have absolutely continuous
representatives on $[0,T]$, with
\begin{equation}
 \begin{aligned}
 K(0)&=\frac12\int\rho_0|u_0|^2\dd,\qquad
 U(0)=c_v\int\rho_0\theta_0\dd,\\
 K'(t)&=-\int Q(u)\dd+\int P\Div u\dd,\\
 U'(t)&=\int Q(u)\dd-\int P\Div u\dd
 \end{aligned}\label{eq:energyidentities}
\end{equation}
for almost every $t$. In particular,
\begin{equation}
 K(t)+U(t)=E_0,\qquad K(t),U(t)\ge0,\qquad 0\le t<T_*.
 \label{eq:energyconservation}
\end{equation}
Moreover, $\int\rho(t)\dd=M$ for all $t<T_*$.
\end{proposition}

\begin{proof}
Fix $T<T_*$. The fixed support, the existing whole-space weak
derivatives, and the finite spatial Lebesgue norms permit the
Poincar\'e inequality on a slightly larger ball. Thus
\begin{equation}
 \begin{aligned}
 &u\in L^\infty(0,T;H^2),\\
 &\theta\in L^\infty(0,T;H^1)\cap L^2(0,T;H^2).
 \end{aligned}\label{eq:compactSobolev}
\end{equation}
Indeed,
$\norm{u}{2}\le C_L\norm{\nabla u}{2}$ and
$\norm{\theta}{2}\le C_L\norm{\nabla\theta}{2}$; the remaining
derivatives are already in \eqref{eq:usedclass}. These are
whole-space Sobolev functions, so the argument introduces no
boundary conditions on an unknown vacuum interface.

The distribution $u_t$ also vanishes outside $\overline{B_L}$ in
spacetime. We claim that
\begin{equation}
 \norm{u_t}{L^2(0,T;L^2)}
       \le C_L\norm{\nabla u_t}{L^2(0,T;L^2)},
 \qquad u_t\in L^2(0,T;H^1).
 \label{eq:utPoincare}
\end{equation}
To prove this for the distributional time derivative, convolve
first in time on a compact subinterval of $(0,T)$ and then in space
with a mollifier of radius $\delta<1$. The resulting smooth function
$w_{\varepsilon,\delta}$ is supported in $B_{L+\delta}$ spatially.
The Poincar\'e inequality on $B_{L+1}$ gives
\[
 \norm{w_{\varepsilon,\delta}(t)}{2}
       \le C_L\norm{\nabla w_{\varepsilon,\delta}(t)}{2}.
\]
Integration on the compact time interval and the $L^2$ contraction
of convolution bound this by the given spacetime norm of $\nabla u_t$.
Weak compactness and distributional convergence identify the limit
with $u_t$. Exhausting $(0,T)$ by compact subintervals proves
\eqref{eq:utPoincare}. The constant is independent of the smoothing
parameters and the compact time subinterval.

Consequently,
$u\in H^1(0,T;H^1)\cap L^\infty(0,T;H^2)$, and its continuous
$H^1$ representative has a limit
\begin{equation}
 u(t)\longrightarrow u_*
       \quad\hbox{in }H^1\quad(t\downarrow0).
 \label{eq:velocitytrace}
\end{equation}
Since $\rho(t)\to\rho_0$ in $W^{1,q}$ and hence in $L^\infty$,
we have $\rho(t)u(t)\to\rho_0u_*$ in $L^2$.
The conservative momentum trace identifies
\begin{equation}
 \rho_0u_*=\rho_0u_0\quad\hbox{almost everywhere}.
 \label{eq:identifytrace}
\end{equation}
This imposes no identification on the initial vacuum set. For the
kinetic energy,
\begin{align}
 &\left|\int\rho(t)|u(t)|^2\dd-\int\rho_0|u_*|^2\dd\right|\notag\\
 &\quad\le
 \norm{\rho(t)-\rho_0}{\infty}\norm{u(t)}{2}^2\notag\\
 &\qquad
   +\norm{\rho_0}{\infty}\norm{u(t)-u_*}{2}
           \bigl(\norm{u(t)}{2}+\norm{u_*}{2}\bigr)
 \longrightarrow0.
 \label{eq:kinetictrace}
\end{align}
Equation \eqref{eq:identifytrace} then gives the asserted value of
$K(0)$.

We derive the balances with integrable right sides. Continuity and
the preceding regularity imply
\begin{equation}
 \rho_t=-u\cdot\nabla\rho-\rho\Div u
       \in L^2(0,T;L^q).
 \label{eq:densitytime}
\end{equation}
Here $u\in L^\infty(0,T;L^\infty)$ follows from
\eqref{eq:compactSobolev}; the other factors use
$\rho,\nabla\rho\in L^\infty(0,T;L^q)$ and
$\Div u\in L^2(0,T;L^\infty)$.
The time product rule for $\int\rho|u|^2$ is legitimate:
for the $\rho_t$ term use
$|u|^2\in L^\infty(0,T;L^{q/(q-1)})$, and for the $u_t$ term
use $\rho\in L^\infty$ and $u,u_t\in L^2$ in space.
These bounds also justify passage from space and time
mollifications to the unmollified products.

Using continuity to write momentum as
$\rho(u_t+u\cdot\nabla u)=\Div S-\nabla P$, we test by $u$,
with a fixed spatial cutoff equal to one near $\overline{B_L}$.
The convective term obeys
\begin{equation}
 \int\rho(u\cdot\nabla u)\cdot u\dd
       =-\frac12\int\Div(\rho u)|u|^2\dd
       =\frac12\int\rho_t|u|^2\dd.
 \label{eq:kinetictransport}
\end{equation}
Weak integration by parts gives
\[
 \int\Div S\cdot u\dd=-\int Q(u)\dd,\qquad
 -\int\nabla P\cdot u\dd=\int P\Div u\dd.
\]
These pairings are valid with $S,P\in L^2$ and $u\in H^1$.
In fact,
\begin{equation}
 \norm{P(t)}{2}
       \le R\norm{\rho(t)}{\infty}^{1/2}
                    \norm{\sqrt{\rho(t)}\theta(t)}{2}.
 \label{eq:pressureLtwo}
\end{equation}
Thus
\begin{equation}
 \frac{d}{dt}\frac12\int\rho|u|^2\dd
       =-\int Q(u)\dd+\int P\Div u\dd.
 \label{eq:kineticbalance}
\end{equation}
The bounds
\begin{equation}
 \begin{aligned}
 \int|Q(u)|\dd&\le C(\mu,\lambda)\norm{\nabla u}{2}^2,\\
 \left|\int P\Div u\dd\right|
       &\le\norm{P}{2}\norm{\Div u}{2}
 \end{aligned}\label{eq:energyintegrability}
\end{equation}
show that the right side is integrable on $[0,T]$. Together with
\eqref{eq:kinetictrace}, this proves the asserted absolute
continuity of $K$.

For internal energy we keep the conservative thermal equation.
Choose $\chi\in C_c^\infty(\R^2)$ equal to one on a neighborhood
of $\overline{B_L}$. Spatial pairing with $\chi$ gives
\begin{equation}
 \frac{d}{dt}\left(c_v\int\rho\theta\,\chi\dd\right)
       =\int Q(u)\dd-\int P\Div u\dd.
 \label{eq:internalbalance}
\end{equation}
Indeed, $\rho\theta u\cdot\nabla\chi=0$. The function $\theta$
and its distributional spatial derivatives are supported in
$\overline{B_L}$ for almost every time, so
\[
 \langle\Delta\theta,\chi\rangle
       =\langle\theta,\Delta\chi\rangle=0.
\]
The functions $Q(u)$ and $P\Div u$ are supported in that ball too.
Thus no flux term at infinity or at a moving vacuum boundary has
been dropped.

The left side of \eqref{eq:internalbalance} is $U'(t)$, and
$U$ is essentially bounded because
$\left|\int\rho\theta\dd\right|
\le\norm{\rho}{1}^{1/2}\norm{\sqrt\rho\,\theta}{2}$.
The right side lies in $L^1(0,T)$ by
\eqref{eq:energyintegrability}, giving an absolutely continuous
extension to $[0,T]$. Pairing the conservative thermal trace
\eqref{eq:traces} with the same $\chi$ yields
\[
 U(0)=c_v\int\rho_0\theta_0\dd.
\]
No unweighted strong trace of $\theta$ is required.
Adding the kinetic and internal balances proves
\eqref{eq:energyconservation}, including its constant at zero.
The signs follow from $\rho,\theta\ge0$ and extend to the
continuous energy representatives. Finally, testing
\eqref{eq:mass} with $\chi$ gives conservation of mass, whose
value is $M$ by the density trace.
\end{proof}

All pointwise-in-time statements about $K$ and $U$ use the absolutely
continuous representatives obtained in Proposition~\ref{prop:energy};
the original Bochner fields are defined only almost everywhere in time.
The equality in \eqref{eq:energyconservation}, with its correctly
identified initial value, is needed below. An upper energy
inequality alone would not give the positive lower bound for the
virial acceleration. The conservative thermal test avoids any
unproved temperature chain rule in vacuum.

\section{The virial bound and small positive-energy data}\label{sec:virial}

\begin{proof}[Proof of Theorem~\ref{thm:obstruction}]
Fix $0<T<T_*$. Let $\chi$ be the cutoff in
Proposition~\ref{prop:energy} and define
\[
 I(t)=\int\rho(t,x)|x|^2\dd,\qquad
 J(t)=\int\rho(t,x)u(t,x)\cdot x\dd.
\]
By fixed support, these are the pairings with $\chi|x|^2$ and
$\chi x$. They are finite; in particular,
$|J(t)|\le L M^{1/2}\norm{\sqrt{\rho(t)}u(t)}{2}$.
Testing continuity with $\chi|x|^2$ gives
\begin{equation}
 I'(t)=2J(t).
 \label{eq:firstvirial}
\end{equation}
Testing conservative momentum with $\chi x$ gives
\begin{equation}
 J'(t)=\int\rho|u|^2\dd-\int\operatorname{tr}S\dd
                                      +2\int P\dd.
 \label{eq:secondvirial}
\end{equation}
Terms involving derivatives of $\chi$ are supported where
$\rho=u=\theta=0$ and their weak spatial derivatives vanish.
In two dimensions,
$\operatorname{tr}S=2(\mu+\lambda)\Div u=2b\Div u$.
Fixed support and $\nabla u\in L^2$ imply $\Div u\in L^1$, and
\[
 \int\Div u\dd=-\int u\cdot\nabla\chi\dd=0.
\]
Consequently,
\begin{equation}
 J'(t)=2K(t)+2\frac{R}{c_v}U(t).
 \label{eq:momentumvirial}
\end{equation}
The right side is integrable, so $J$ has an absolutely continuous
representative. The conservative momentum trace tested with $\chi x$
gives $J(0)=J_0$; the density trace gives $I(0)=I_0$.
Equation \eqref{eq:firstvirial} identifies the absolutely continuous
derivative of $I$. Nonnegativity and conservation of energy give
\begin{equation}
 I''(t)=4K(t)+4\frac{R}{c_v}U(t)
       \ge4\beta\bigl(K(t)+U(t)\bigr)=4\beta E_0
 \label{eq:virialacceleration}
\end{equation}
almost everywhere. Integrating twice,
\begin{equation}
 I(t)\ge I_0+2J_0t+2\beta E_0t^2,\qquad 0\le t\le T.
 \label{eq:quadraticgrowth}
\end{equation}
On the other hand, fixed support and mass conservation yield
\begin{equation}
 0\le I(t)\le L^2\int\rho(t)\dd=ML^2.
 \label{eq:momentupper}
\end{equation}
This proves \eqref{eq:mainbound}. Since $E_0>0$ and
$ML^2-I_0>0$, the inequality
\[
 2\beta E_0t^2+2J_0t-(ML^2-I_0)\le0
\]
has exactly one positive root, the right side of \eqref{eq:lifespan}.
Every $t<T_*$ is bounded by it, since $T<T_*$ was arbitrary.
The theorem follows.
\end{proof}

We finish with smooth initial data whose common norm bounds can be
fixed before the energy parameter is chosen.

\begin{corollary}\label{cor:smallenergy}
Fix the parameters in \eqref{eq:parameters}, $a>1$, $q>2$,
$\eta_0>0$, and a mass $M>0$. There is a family
$(\rho_0^\varepsilon,u_0^\varepsilon,\theta_0^\varepsilon)$,
$0<\varepsilon\le1$, of smooth compactly supported data satisfying
\eqref{eq:data}--\eqref{eq:compatibility} such that
\begin{equation}
 \begin{gathered}
 \int\rho_0^\varepsilon\dd=M,\qquad
 \supp\rho_0^\varepsilon\subset\overline{B_1},\\
 E_0^\varepsilon>0,\qquad E_0^\varepsilon\longrightarrow0.
 \end{gathered}
 \label{eq:familyproperties}
\end{equation}
The density is independent of $\varepsilon$. Every norm displayed
in \eqref{eq:data}, as well as $\norm{g^\varepsilon}{2}$, has a
common finite upper bound. No member of the family admits a global
solution in Definition~\ref{def:solution}.
\end{corollary}

\begin{proof}
Take the nonnegative smooth function
\[
 f(x)=
 \begin{cases}
 \exp\bigl(-1/(1-|x|^2)\bigr),&|x|<1,\\
 0,&|x|\ge1,
 \end{cases}
 \qquad
 c=\frac{M}{\int f^2\dd},
 \qquad r=cf^2.
\]
Then $r,\sqrt r\in C_c^\infty(\R^2)$, $r>0$ on $B_1$, and
$\int r\dd=M$. Set
\begin{equation}
 \rho_0^\varepsilon=r,\qquad
 u_0^\varepsilon=0,\qquad
 \theta_0^\varepsilon=\varepsilon r,
 \qquad 0<\varepsilon\le1.
 \label{eq:familydefinition}
\end{equation}
The weighted density norm
$\norm{\bar x^a r}{L^1\cap H^1\cap W^{1,q}}$ is fixed and finite.
The velocity norms vanish, and
\begin{equation}
 \begin{aligned}
 \norm{\sqrt r\,\theta_0^\varepsilon}{2}
       &=\varepsilon\left(\int r^3\dd\right)^{1/2},\\
 \norm{\nabla\theta_0^\varepsilon}{2}
       &=\varepsilon\norm{\nabla r}{2}.
 \end{aligned}\label{eq:familynorms}
\end{equation}
The momentum compatibility condition holds exactly:
\begin{align}
 &-\mu\Delta u_0^\varepsilon
       -b\nabla\Div u_0^\varepsilon
       +R\nabla(\rho_0^\varepsilon\theta_0^\varepsilon)\notag\\
 &\hspace{16mm}=R\varepsilon\nabla(r^2)
       =2R\varepsilon r\nabla r
       =\sqrt r\,g^\varepsilon,\notag\\
 &g^\varepsilon=2R\varepsilon\sqrt r\,\nabla r.
 \label{eq:familycompatibility}
\end{align}
Since $\sqrt r\in C_c^\infty$, $g^\varepsilon$ is smooth and
compactly supported, with
\begin{equation}
 \norm{g^\varepsilon}{2}
       =2R\varepsilon
                \left(\int r|\nabla r|^2\dd\right)^{1/2}.
 \label{eq:familygbound}
\end{equation}
All required initial bounds are therefore uniform for
$0<\varepsilon\le1$. No thermal compatibility condition is imposed.

Finally,
\begin{equation}
 E_0^\varepsilon=c_v\varepsilon\int r^2\dd>0,\qquad
 J_0^\varepsilon=0,\qquad
 I_0^\varepsilon=\int r|x|^2\dd<M.
 \label{eq:familyenergy}
\end{equation}
Theorem~\ref{thm:obstruction} applies to each $\varepsilon>0$.
If a solution in the specified class exists on $[0,T_*^\varepsilon)$,
its interval satisfies
\begin{equation}
 T_*^\varepsilon
       \le\left(\frac{M-\int r|x|^2\dd}
                 {2\beta c_v\varepsilon\int r^2\dd}\right)^{1/2}.
 \label{eq:familylifespan}
\end{equation}
This is finite for every $\varepsilon>0$.
\end{proof}

The family has fixed positive mass, fixed density support, and zero
total momentum. The initial velocity and temperature are actual
compactly supported smooth functions, so their integrability has no
constant-representative ambiguity. A positive energy threshold based
only on the fixed parameters and these common initial upper bounds
therefore cannot guarantee a global solution in the stated class for
every member of the family. The upper bound
\eqref{eq:familylifespan} tends to infinity as
$\varepsilon\downarrow0$; the conclusion concerns each strictly
positive energy.

The construction verifies the displayed initial assumptions and
compatibility condition, but does not independently establish local
solvability in \eqref{eq:fullclass}. It gives no example of finite-time
breakdown without that additional result. The theorem also retains
nonnegative temperature, finite unweighted velocity and temperature
norms, and strict $\mu+\lambda>0$. Different solution classes and the
viscosity endpoint require separate arguments.

\begin{thebibliography}{1}
\bibitem{WangV3}
Wang, X.:
On Existence and Blowup Criterion of Strong Solutions to
Cauchy Problem of 2D Full Compressible Navier--Stokes System with Vacuum.
arXiv:2212.13343v3 [math.AP], 4 April 2023.
\end{thebibliography}
\end{document}
